\clearreportauthor
\section{Conclusion}\label{sec:conclusion}

This project achieved its objective of designing, implementing, and evaluating a reproducible pipeline for improving the routing of a multi-agent chatbot. The resulting system integrates a selective similarity-based Stage 1 Router, an LLM-based Stage 2 Router, an LLM judge, and an OPRO-based prompt optimizer in a working showcase application. The accompanying dashboards make the evaluation, training, optimization, and deployment workflow inspectable rather than leaving prompt improvement as an ad hoc manual process.

The controlled evaluation provides evidence that the approach can improve the agreed performance dimensions. On the fixed 240-query evaluation dataset, optimizing the Stage 2 prompt increased routing accuracy from 83\% to 100\%. The complete two-stage Router achieved 98\% accuracy, routed 78\% of queries without a Stage 2 LLM call, and reduced average latency from 669 ms to 407 ms. It therefore retained most of the accuracy gain while lowering model use and latency. At the same time, the difference between judge-based and golden-label results shows that automated feedback is useful for optimization but should not be treated as perfectly calibrated ground truth.

The findings remain subject to the project's main limitations: all experiments used synthetic data, the improvement loop relied on a single judge, and only a limited set of models and configurations was tested. The results therefore demonstrate technical feasibility in the showcase environment, not production performance at Interhyp. A next step should validate the pipeline on independently labelled production samples, monitor accuracy, latency, and cost under live traffic, and introduce periodic human review or a judge panel if reliability declines. With these safeguards, the implemented pipeline offers Interhyp a practical foundation for replacing case-by-case Router refinement with a systematic and measurable improvement process.
